\section{Identifying stages of soot formation}
\label{sec:paper3_stages_soot_formations}

This section presents the paper titled ``\textit{Yield, morphology and maturity of carbon black made by methane pyrolysis in a shock tube}'' and submitted to \textit{Carbon journal}, and it is focused on using Omnisoot to identify different stages of soot formation during methane pyrolysis in a shock-tube.

Despite extensive studies on the chemical kinetics of methane pyrolysis (Refs. in \cite{fau2013methane}), the mechanisms governing soot formation remain poorly understood. This limitation arises from incomplete knowledge of the pathways leading to soot inception and from uncertainties in existing kinetic mechanisms~\citep{agafonov2016unified}. These uncertainties originate with methane dissociation reactions~\citep{vlasov2022experimental} and propagate through key intermediates such as methyl ($\mathrm{CH_3}$)~\citep{wang2016improved}, ethylene ($\mathrm{C_2H_4}$), and acetylene ($\mathrm{C_2H_2}$)~\citep{sajid2016shock}. The uncertainty is further amplified for polycyclic aromatic hydrocarbons (PAHs)~\citep{nativel2019shock}, which are widely accepted as the primary precursors to soot~\cite{alfe2009structure}. Previous studies have shown that predicted concentrations of large PAHs can vary by orders of magnitude depending on the kinetic mechanism employed~\citep{wang2023systematic}. This variability underscores the need for time-resolved measurements of gas-phase species and particulate properties under well-controlled conditions to improve quantitative understanding of methane conversion, intermediate formation, and carbon flux to soot through inception and surface growth.

Shock tubes are powerful tools for this purpose because a wide range of temperatures (1700-3000~K~\citep{agafonov2016unified}) and pressures (1-100~atm~\citep{shao2020shock}) can be achieved by adjusting the pre-shock pressure ratio; the instant heating of the mixture behind the reflected shock creates a nearly isothermal and isobaric condition over a short test time (1-3 ms) with minimal diffusion and mixing~\citep{eremin2012formation}, which is an ideal setting for kinetic studies of pyrolysis pathways and soot inception and surface growth.

Capturing the rapid decomposition of methane and the subsequent formation of intermediates can be achieved with the high temporal resolution offered by continuous-wave (CW) laser absorption spectroscopy~\citep{pinkowski2019multi}. In this technique, the attenuation of light in an absorbing medium is related to the gas density through the Beer–Lambert law, provided that the absorption cross-section is known. Accurate application, however, is complicated by the strong dependence of the absorption cross-section on temperature, pressure, and gas composition. soot yield and volume fraction can also be quantified with good accuracy from light extinction measurements based on the same principle, as long as the particles remain in the Rayleigh limit and the particle volume fraction is below 10~ppm~\citep{eremin2012formation}. This approach requires prior knowledge of the optical properties of soot, represented by the absorption function $E(m)$. 

While a constant $E(m)$ value of mature soot is commonly used to retrieve volume fractions from extinction data~\citep{de2008scattering, agafonov2016unified, utsav2017simultaneous}, the optical properties are noticeably different for incipient and mature soot~\citep{michelsen2021effects}. During maturation, soot particles undergo structural and chemical changes that result in the reduction of hydrogen-to-carbon ratio (H/C) due to carbonization~\cite{kholghy2016core} that increases the density and specific heat~\citep{michelsen2021effects}. The maturity level of soot is linked to the spectral dependence of absorption~\citep{bescond2016soot} quantified by the ratio of $E(m)$ at two different wavelengths~\citep{yon2021revealing} that approaches unity for mature soot; therefore, performing extinction measurements at multiple wavelengths can reduce the uncertainties of $E(m)$ and provide a good measure of soot maturity.


This section deals with simultaneous measurements of gaseous species and soot during the pyrolysis of 5\% $\mathrm{CH_4}$ at post-reflected-shock pressures near 5~atm and temperatures between 1800 and 2500~K. In that work, a computational framework coupling detailed gas-phase chemistry with soot inception, surface growth, and coagulation is developed to describe methane-to-soot conversion, successfully reproducing the measured time histories of key gas-phase species and soot volume fraction, as well as the final primary particle diameter. A detailed analysis of the simulation results, complemented by experimental measurements, is conducted to delineate distinct stages of soot formation based on the time evolution of $f_v$ and carbon yield. Within each stage, the evolution of inception flux and surface growth mechanisms is examined to elucidate their influence on $f_v$, yield, and morphology, as well as their coupling with gas-phase chemistry. The deviation of model predictions from $f_v$, $d_p$, and induction time measurements are discussed in detail. Particular emphasis is placed on the role of carbonization in describing soot maturity, its effect in limiting surface growth rates, and its connection to the spectral dependence of soot particles characterized through multi-wavelength light extinction.